\section{Ausgangssituation der Informationsarchitektur und UX}
\label{sec:03-03_ia-and-ux-current-state}

Die Voranalyse zeigte, dass die \Gls{ia} primär historisch gewachsen und stark von internen Organisationslogiken geprägt war. Inhalte waren häufig entlang von Projekten, Zuständigkeiten oder internen Bezeichnungen strukturiert, statt sich an Aufgaben und mentalen Modellen der Nutzenden zu orientieren \cite{419:Glushko2013}. Für neue oder sporadische Nutzende führte dies zu Orientierungsproblemen, da Begriffe, Kategorien und Navigationspfade ohne explizites Domänenwissen schwer interpretierbar waren. Zentrale Informationen waren zwar vorhanden, jedoch nicht konsistent auffindbar oder eindeutig verortet; dieselben Inhalte traten teilweise in unterschiedlichen Kontexten mit abweichenden Bezeichnungen auf, was die Findbarkeit zusätzlich verringerte \cite{419:Glushko2013,307:DWGFindability}.

Aus UX-Sicht äußerten sich diese strukturellen Defizite vor allem in erhöhten Suchzeiten, unsicherem Navigationsverhalten und einer starken Abhängigkeit von individueller Erfahrung mit der Plattform. \Gls{usability}-Tests im Ist-Zustand verdeutlichten, dass selbst grundlegende Aufgaben -- etwa das Finden aktueller Inhalte, das Nachschlagen von Begriffen oder das Wiederfinden relevanter Seiten -- unnötigen kognitiven Aufwand erforderten, weil Nutzende wiederholt zwischen Navigation und Suche wechseln, mehrere Einstiege ausprobieren oder sich auf inoffizielle Linksammlungen verlassen mussten \cite{113:NNGIntranetGuidelines}. Nielsen-Norman-Analysen zu Intranet-\Gls{usability} zeigen ähnliche Muster: Schwächen in Informationsarchitektur und Sucherfahrung gehören zu den häufigsten Ursachen für Produktivitätsverluste, und die Qualität von Navigation und Suche erklärt einen erheblichen Anteil der Leistungsunterschiede zwischen gut und schlecht nutzbaren Intranets \cite{113:NNGIntranetGuidelines,111:NNGSearchVisible}. Die im Projekt beobachteten Probleme lassen sich somit als typisches Beispiel eines historisch gewachsenen Enterprise-Wikis einordnen, in dem fehlende Governance und uneinheitliche Strukturierungsprinzipien die effektive Nutzung der Plattform begrenzen \cite{307:DWGFindability}.
